\section{Introduction}
\label{sec:intro}

Behavior-driven development expresses expected software behavior in executable or executable-adjacent scenarios, commonly using Gherkin. Its cost is not limited to writing Given--When--Then syntax: engineers must preserve domain preconditions, project conventions, data dependencies and expected outcomes while keeping feature files parsable by the project toolchain. In the three projects of this study, a Sylius checkout scenario may depend on channel, country, shipping-method state and confirmation behavior, Apache Fineract scenarios encode dates, repayment periods, balances and data tables, and Diaspora scenarios describe browser-visible behavior such as navigation, tags, comments and notifications. A generated scenario can therefore be syntactically valid yet incomplete as a specification of business behavior, and converting a compact Connextra story into project-consistent Gherkin is a constrained generation problem rather than a simple paraphrasing task.

\subsection{Research Gap and Motivation}

Prior work shows that LLMs can generate BDD artifacts~\cite{rathnayake2026,karpurapu2024,ferreira2025,fernandes2025}, but the evaluations differ substantially in data scale, model family and validation method (Section~\ref{sec:related}). These studies establish feasibility, but they do not jointly answer whether a current compact model can produce both parser-ready and structurally close Gherkin across multiple public software domains. In particular, none of the five studies in Table~S5 of the supplementary material tests whether its semantic measure reacts to business-logic defects, and none varies the example given in the prompt.

This study addresses that narrower gap by evaluating the pinned gpt-5.4-mini-2026-03-17 snapshot on 261 stories from three unrelated open-source projects, using same-project one-shot prompting and two separate outcome measures. The contribution is not a claim that GPT-5.4-mini is superior to earlier models, but an auditable evaluation design that separates syntactic validity from structural similarity on a public, multi-domain benchmark and tests how far the similarity metric can be trusted. Because the stories were reverse-engineered from the reference files and the 0.85 reference point was not calibrated, this study therefore includes a leakage measurement, a defect audit, a perturbation test, a reference-point sensitivity analysis, per-project results, a prompt ablation and a same-time rerun.

\subsection{Research Questions and Contributions}

\textbf{RQ1 -- Structural skeleton resemblance.} When GPT-5.4-mini reconstructs feature files from user stories reverse-engineered from them, how is the skeleton cosine similarity to the reference distributed around the reference point of 0.85 fixed in the proposal?

\noindent\textbf{RQ2 -- Syntactic validity.} Does GPT-5.4-mini generate feature files that the behave Gherkin parser accepts at a rate above 80\%?

The paper makes three contributions: (1) a 261-case, multi-domain benchmark derived from project-maintained reference files in Sylius, Apache Fineract, and Diaspora, with its construction and leakage risk documented; (2) an evaluation that reports parser compliance separately from structural similarity and checks the similarity metric against defect labels, seeded defects and alternative reference points; and (3) a replication package containing the pinned model identifier, prompts, scripts, API logs, outputs, and all additional analyses.

Section~\ref{sec:related} reviews related work. Section~\ref{sec:method} describes the dataset, the generation pipeline, the metrics and the analyses, and Section~\ref{sec:results} reports the results. Sections~\ref{sec:discussion} and~\ref{sec:threats} discuss the implications and the threats to validity, and Section~\ref{sec:conclusion} concludes. The supplementary material, linked under Data and Ethics, gives the prompts and additional results (Section~S1), the perturbation test (Section~S2) and further method and result details (Section~S3).
